% ============================================================
% Plantilla base — Reportes en formato Springer LNCS (llncs)
% Copia este fichero como main.tex dentro de la carpeta del reporte
% (reportes/datapath-1/main.tex, etc.) y rellena el contenido.
%
% Compilar SIEMPRE con:
%   ../../scripts/latexbuild main.tex      (desde la carpeta del reporte)
%   → genera build/main.pdf
%
% REQUISITOS DE EXTENSIÓN (ver reportes/AGENTS.md):
%   Resumen < 100 palabras | Introducción 1-2 cuartillas
%   Diseño teórico 3-5 | Propuesta experimental 1 hoja
%   Simulaciones (capturas) máx 2 hojas | Resultados 3-5
%   Conclusiones 1 | Referencias >= 15 | Apéndice: código VHDL
% ============================================================
\documentclass[llncs]{llncs}

\usepackage[utf8]{inputenc}
\usepackage[T1]{fontenc}
\usepackage{graphicx}
\usepackage{listings}
\usepackage{hyperref}
\usepackage{xcolor}

% ---- Estilo para código VHDL (apéndice y fragmentos) ----
\lstdefinestyle{vhdl}{
  language=VHDL,
  basicstyle=\ttfamily\scriptsize,
  keywordstyle=\color{black},
  commentstyle=\color{gray}\itshape,
  stringstyle=\color{black},
  showstringspaces=false,
  breaklines=true,
  frame=single,
  rulecolor=\color{black}
}
\lstset{style=vhdl}

\begin{document}

% ============================================================
% TÍTULO Y AUTORES
% ============================================================
\title{TÍTULO DEL REPORTE}

\author{Nombre Apellido$^{1}$,%
\and Nombre Apellido$^{1}$,%
\and Nombre Apellido$^{1}$,%
\and Nombre Apellido$^{1}$,%
\and Nombre Apellido$^{1}$}

\authorrunning{Apellido et al.}
\institute{Equipo de Arquitectura de Computadoras,\\
Universidad, Ciudad, País\\
\email{\{correo1, correo2, correo3, correo4, correo5\}@universidad.mx}}

\maketitle

% ============================================================
% RESUMEN  (< 100 palabras — fuente 9pt con sangría, estilo llncs)
% ============================================================
\begin{abstract}
Resumen. Escribe aquí la síntesis del trabajo en \textbf{menos de 100
palabras}: objetivo, metodología (GHDL/GTKWave, VHDL-2008) y resultado
principal.
\end{abstract}

\noindent\textbf{Palabras Clave:} Arquitectura de Computadoras, Datapath,
VHDL, Simulación, GHDL

% ============================================================
% 1. INTRODUCCIÓN  (1–2 cuartillas)
% ============================================================
\section{Introducción}
\label{sec:intro}

Contexto, motivación y objetivo del trabajo. Indica la estructura del
artículo. (1--2 cuartillas.)

% ============================================================
% 2. DISEÑO TEÓRICO  (3–5 cuartillas)
% ============================================================
\section{Diseño Teórico}
\label{sec:diseno}

Marco teórico, arquitectura del datapath, descripción de los módulos VHDL
entidades, registros, ALU, buses) y decisiones de diseño. (3--5 cuartillas.)

% ============================================================
% 3. PROPUESTA EXPERIMENTAL  (1 hoja)
% ============================================================
\section{Propuesta Experimental}
\label{sec:propuesta}

Metodología: herramientas de simulación (GHDL~\cite{ghdl-docs},
GTKWave~\cite{gtkwave-docs}), estándar VHDL-2008, bancos de
prueba, variables y procedimiento de evaluación. (1 hoja.)

% ============================================================
% 4. SIMULACIONES DE VHDL  (capturas de los módulos más importantes)
% ============================================================
\section{Simulaciones de VHDL}
\label{sec:simulaciones}

Capturas de los módulos más importantes (máximo 2 hojas). Exporta las ondas
de \texttt{vhdl/<unidad>/salidas/*.vcd} a PNG en \texttt{figuras/}.

% Descomenta el bloque cuando exista la captura en figuras/
% \begin{figure}[t]
%   \centering
%   \includegraphics[width=0.9\textwidth]{figuras/captura_modulo.png}
%   \caption{Simulación de <módulo>: descripción breve de lo que se observa.}
%   \label{fig:captura}
% \end{figure}

% ============================================================
% 5. RESULTADOS  (3–5 cuartillas)
% ============================================================
\section{Resultados}
\label{sec:resultados}

Análisis de las simulaciones: tablas comparativas, tiempos, comportamiento
observado, cumplimiento de las pruebas (assert). Interpretación crítica.
(3--5 cuartillas.)

% ============================================================
% 6. CONCLUSIONES  (1 cuartilla)
% ============================================================
\section{Conclusiones}
\label{sec:conclusiones}

Cierre: logros frente a los objetivos, limitaciones y trabajo futuro.
(1 cuartilla.)

% ============================================================
% REFERENCIAS  (>= 15 — indicar exactamente dónde aparece el texto)
% ============================================================
% IMPORTANTE: cada referencia debe permitir ubicar exactamente dónde
% aparece el texto sugerido (capítulo/§/página, o URL + fecha de consulta).
% Anótala en el campo note de la entrada (más seguro que un comentario,
% porque el .bib no admite comentarios dentro de una entrada).
% Crea también referencias.bib junto a main.tex (hay ejemplos en
% reportes/<unidad>/referencias.bib). Ojo: bibtex falla si no hay
% ninguna \cite en el texto.
\bibliographystyle{splncs04}
\bibliography{referencias}

% ============================================================
% APÉNDICE — CÓDIGO VHDL
% ============================================================
\appendix
\section{Código VHDL}
\label{app:codigo}

Código completo de los módulos de la unidad (copiado de
\texttt{vhdl/<unidad>/src/}). Usa \texttt{lstlisting}:

\begin{lstlisting}[caption={sumador.vhd --- <descripción>}, label=lst:sumador]
library ieee;
use ieee.std_logic_1164.all;
use ieee.numeric_std.all;

entity sumador is
  port (
    a, b : in  unsigned(3 downto 0);
    s    : out unsigned(3 downto 0)
  );
end entity;

architecture rtl of sumador is
begin
  s <= a + b;
end architecture;
\end{lstlisting}

\end{document}
